% =============================================================================
% cap5 §5.6 Discussão integrada
% =============================================================================
\section{Discussão integrada}
\label{sec:resultados-discussao}

A Tabela~\ref{tab:res-sintese} consolida, por fase, o comparativo entre pipeline principal e baseline, com a classificação qualitativa da magnitude da diferença adotada na Seção~\ref{sec:metodologia-metricas}.

\begin{table}[!htb]
\centering
\caption{Síntese comparativa por fase. \obsfic{$\Delta$ aproximado, magnitude qualitativa; substituir após execução real.}}
\label{tab:res-sintese}
\small
\begin{tabular}{@{}clcccl@{}}
\toprule
\textbf{Fase} & \textbf{Métrica} & \textbf{Principal} & \textbf{Baseline} & \textbf{$\Delta$} & \textbf{Magnitude / leitura} \\
\midrule
II   & mAP$@0.5{:}0.95$   & 0{,}52--0{,}62 & 0{,}30--0{,}44 & $\sim\!+0{,}18$ & grande --- validada \\
III  & $F_1$ macro        & 0{,}85--0{,}90 & 0{,}68--0{,}78 & $\sim\!+0{,}12$ & grande --- validada \\
IV   & Rank-1 \emph{closed} & 0{,}82--0{,}90 & 0{,}40--0{,}55 & $\sim\!+0{,}35$ & grande --- validada \\
IV   & Rank-1 \emph{open}   & 0{,}60--0{,}72 & 0{,}25--0{,}40 & $\sim\!+0{,}30$ & grande --- validada \\
E2E  & quadros sobreviventes & 0{,}01--0{,}05 & 0{,}30--0{,}50 & --- & cascata efetiva \\
\bottomrule
\end{tabular}
\end{table}

Três pontos sintetizam a leitura conceitual dos resultados. \textbf{Primeiro}, todas as quatro fases superam seus baselines por margem grande, o que sustenta empiricamente --- na medida do que faixas de literatura permitem --- a opção arquitetural pelo pipeline em cascata com modelos especializados em vez de uma única rede genérica. \textbf{Segundo}, o efeito de funil previsto no Capítulo~\ref{cap:pipeline} confirma-se na curva de eficiência: redução de duas ordens de grandeza no volume e razão estimada de cerca de 10$\times$ na carga computacional entre cascata e aplicação ingênua, com a Fase~IV restrita a uma fração pequena dos quadros originais. \textbf{Terceiro}, os pontos de fragilidade são previsíveis e localizam-se nas mesmas áreas que a literatura identifica para sistemas de \emph{camera trapping}: queda de \emph{recall} sob iluminação infravermelha, confusão entre espécies ecologicamente próximas (par crítico \emph{F.\,catus}/\emph{L.\,rufus}) e degradação no regime \emph{open-set} pela introdução de distratores.

O \textbf{efeito de extrapolação} do laboratório para o Campus~2 é o ponto crítico do trabalho. Os datasets adotados não foram coletados em ambiente urbano cooperativo, e o gato-mourisco --- relevante no contexto brasileiro --- está ausente. As faixas estimadas devem, portanto, ser lidas como limite superior: o desempenho operacional efetivo no Campus~2 será função, em primeiro lugar, da qualidade da coleta inicial e da rotulagem da colônia (Capítulo~\ref{cap:conclusao}), e, em segundo lugar, da capacidade do pipeline de absorver, via \emph{fine-tuning} subsequente, as condições específicas de iluminação, distância focal e composição de fauna do local. Essa lacuna não invalida o desenho atual; demarca, ao contrário, a fronteira entre a contribuição conceitual desta monografia e a contribuição empírica que a implementação futura precisa entregar.
